\section{Related Work} \label{sec:related-work}

This section explores current research in \ac{llms} and their application to \ac{aqg}, with sections to detail the literature review process and highlight significant contributions in the field.

% \newtodo{Discuss other work in the field aiming this related research problem.}
% \newtodo{Show how my research builds on them, where used for inspiration, and e.g. how approaches may be combined.}
% \newtodo{Mention papers, explain why good or bad, also in comparison with my approach.}
% \newtodo{Try to categorize and group the related work by these, show (dis)advantages using one example.}

\object{Literature Review} To find the fundamental principles and state-of-the-art of the research problem, a literature review was conducted. It was performed in the following way: First, several searches were done with respect to \ac{llms} in common, question generation, particularily focusing on Bloom's Taxonomy, and lastly evaluation methods for both categories. Google Scholar and Elicit were used to find relevant papers. Several search terms were coupled in Google Scholar, such as \textit{LLM, Large Language Model, architecture, question generation, automatic question generation, Bloom's Taxonomy, evaluation, assessment, alignment}. Elicit is a search engine that is designed to help researchers find relevant papers based on writing a specific prompt, instead of a query of keywords.

\object{Review Steps} The review steps were as follows: 
\begin{itemize}
    \item First, the papers were downloaded and stored in Zotero. Then, duplicates were filtered out.
    \item The remaining papers were \textbf{included} or excluded by using the TACID method. Papers were included if they are recent (published from 2021 on) and dealing with \ac{llms} or automated question generation or evaluation methods. If the paper was a preprint instead of a peer-reviewed paper, it needed a comparatively high relevance to be included ($\approx 10$ citations). If a publication was present in a workshop, it needed citations to be included.
    \item Papers were \textbf{excluded} if they were not written in English, not open access, or the journal is predatory\footnote{Journals noted on \url{https://beallslist.net/}, accessed 04/30/2025}. If the focus is explicitly on \ac{slms} or multimodal models, or other domains besides the desired ones, the paper was excluded. 
    \item After the process, the remaining papers were filtered by a deeper analysis of the papers.
    \item In the end, diverse cited and subsequently found papers blog posts by e.g.\ Google were added to the thesis to improve the information space.
\end{itemize}

\subsection{Current LLMs Research}

\hspace{1.5em}\object{Review Challenges at \ac{llms}} Finding the current state-of-the-art of \ac{llms} with the explained review method did not yield recent results, even if papers such as surveys were published in 2024 or 2025. The reason for this is that the field of \ac{llms} is moving so fast, that the review papers become outdated. Moreover, the trend of the models and their releases were followed manually, as captured in Table \bfref{tab:llm_releases} below. It shows the most important releases of \ac{llms} in the past years since ChatGPT \cite{openai_introducing_2022}. The table is not complete, but it shows well-known models and their architecture, including \textit{reasoning} capabilities, if these are known to the public. \textit{Reasoning} means, in the OpenAI o1 system card \cite{openai_openai_2024-1} for instance, the model thinks in a manner of \ac{cot} before returning the desired response to the user. The proposed table is sorted by the release date of the models to obtain a timeline and a way to interpret the development of the models.

\begin{table}[htbp]
    \centering
    \caption{Timeline of Major LLM Releases (2022-2025).}
    \label{tab:llm_releases}
    \renewcommand{\arraystretch}{1.4}
    \definecolor{openaiColor}{HTML}{c6e2ce}   
    \definecolor{metaColor}{HTML}{9deade}     
    \definecolor{googleColor}{HTML}{ebd9a1}   
    \definecolor{anthropicColor}{HTML}{cdd8f2}
    \definecolor{xaiColor}{HTML}{efd4d0} 
    \definecolor{deepseekColor}{HTML}{eed6b3} 
    \definecolor{alibabaColor}{HTML}{cde3b5}
    % https://medialab.github.io/iwanthue/

    \begin{tabularx}{\textwidth}{|>{\bfseries\centering\arraybackslash}p{2.5cm}|>{\centering\arraybackslash}p{2.25cm}|>{\centering\arraybackslash}p{4.5cm}|>{\centering\arraybackslash}p{3cm}|>{\centering\arraybackslash}X|}
        \hline
        \rowcolor{gray!15} \textbf{Date} & \textbf{Developers}                            & \textbf{Model Releases} & \textbf{Architecture} & \textbf{Reasoning} \\
        \hline
        Nov 2022                         & \cellcolor{openaiColor}{\textbf{OpenAI}}       & ChatGPT/GPT-3.5         & Decoder-only?         & -- \\ % https://openai.com/index/chatgpt/
        \hline
        \multirow{2}{*}{Feb 2023}        & \cellcolor{metaColor}{\textbf{Meta}}           & LLaMA 65B               & Transformer           & -- \\ % https://arxiv.org/pdf/2302.13971 paper (not really known architecture)
        \cline{2-5}
                                         & \cellcolor{googleColor}{\textbf{Google}}       & Bard 137B               & Decoder-only          & -- \\ % https://ai.google/advancing-ai/milestones/?section=bard is decoder only cuz powered by lamda 137B https://arxiv.org/pdf/2201.08239
        \hline
        \multirow{2}{*}{Mar 2023}        & \cellcolor{openaiColor}{\textbf{OpenAI}}       & GPT-4                   & Decoder-only?         & -- \\ % https://arxiv.org/pdf/2303.08774
        \cline{2-5}
                                         & \cellcolor{anthropicColor}{\textbf{Anthropic}} & Claude 1                & --                     & -- \\ % https://www.anthropic.com/news/introducing-claude
        \hline
        \multirow{2}{*}{Jul 2023}        & \cellcolor{anthropicColor}{\textbf{Anthropic}} & Claude 2                & --                     & -- \\ % https://www.anthropic.com/news/claude-2
        \cline{2-5}
                                         & \cellcolor{metaColor}{\textbf{Meta}}           & LLaMA 2 70B             & Decoder-only          & -- \\ % https://arxiv.org/pdf/2307.09288 paper
        \hline
        Nov 2023                         & \cellcolor{xaiColor}{\textbf{xAI}}             & Grok-1 314B             & MoE                   & -- \\ % https://x.ai/news/grok & https://x.ai/news/grok-os (here MoE)
        \hline
        Dec 2023                         & \cellcolor{googleColor}{\textbf{Google}}       & Gemini                  & Decoder-only          & -- \\ % https://arxiv.org/pdf/2312.11805
        \hline
        Feb 2024                         & \cellcolor{googleColor}{\textbf{Google}}       & Gemini 1.5              & Decoder MoE           & -- \\ % https://arxiv.org/pdf/2403.05530
        \hline
        Mar 2024                         & \cellcolor{anthropicColor}{\textbf{Anthropic}} & Claude 3                & --                     & -- \\ % https://assets.anthropic.com/m/61e7d27f8c8f5919/original/Claude-3-Model-Card.pdf
        \hline
        % Apr 2024                         & \cellcolor{metaColor}{\textbf{Meta}}           & LLaMA 3 70B             & Decoder-only          & -- \\ % https://ai.meta.com/blog/meta-llama-3/
        % \hline
        May 2024                         & \cellcolor{openaiColor}{\textbf{OpenAI}}       & GPT-4o                  & Decoder-only?         & -- \\ % https://arxiv.org/pdf/2410.21276
        \hline
        Jun 2024                         & \cellcolor{anthropicColor}{\textbf{Anthropic}} & Claude 3.5 Sonnet       & --                     & -- \\ % https://www.anthropic.com/news/claude-3-5-sonnet
        \hline
        % Jul 2024                         & \cellcolor{metaColor}{\textbf{Meta}}           & LLaMA 3.1 405B          & Decoder-only          & -- \\ % https://arxiv.org/pdf/2407.21783 paper (dense decoder transformer instead of MoE, presented in Shazeer et al. 2017)
        % \hline
        Aug 2024                         & \cellcolor{xaiColor}{\textbf{xAI}}             & Grok-2                  & --                     & -- \\ % https://x.ai/news/grok-2
        \hline
        Sep 2024                         & \cellcolor{openaiColor}{\textbf{OpenAI}}       & o1-mini                 & --                     & Yes \\ % https://openai.com/index/openai-o1-mini-advancing-cost-efficient-reasoning/
        \hline
        \multirow{4}{*}{Dec 2024}        & \cellcolor{metaColor}{\textbf{Meta}}           & LLaMA 3.3 70B           & Decoder-only          & -- \\ % https://ai.meta.com/blog/future-of-ai-built-with-llama/ and https://huggingface.co/meta-llama/Llama-3.3-70B-Instruct (decoder-only there)
        \cline{2-5}
                                         & \cellcolor{googleColor}{\textbf{Google}}       & Gemini 2.0              & --                     & -- \\ % https://blog.google/technology/google-deepmind/google-gemini-ai-update-december-2024/
        \cline{2-5}
                                         & \cellcolor{openaiColor}{\textbf{OpenAI}}       & o1/o3-mini              & --                     & Yes \\ % https://arxiv.org/pdf/2412.16720 o1 paper, https://openai.com/index/openai-o3-mini/
        \cline{2-5}
                                         & \cellcolor{deepseekColor}{\textbf{DeepSeek}}   & DeepSeek-V3 658B        & MoE                   & -- \\ % https://arxiv.org/pdf/2412.19437
        \hline
        \multirow{2}{*}{Jan 2025}        & \cellcolor{deepseekColor}{\textbf{DeepSeek}}   & DeepSeek R1 658B        & MoE                   & Yes \\ % https://arxiv.org/pdf/2501.12948
        \cline{2-5}
                                         & \cellcolor{alibabaColor}{\textbf{Alibaba}}     & Qwen 2.5-Max            & MoE                   & -- \\ % https://qwenlm.github.io/blog/qwen2.5-max/
        \hline
        \multirow{3}{*}{Feb 2025}        & \cellcolor{xaiColor}{\textbf{xAI}}             & Grok-3                  & --                     & Yes \\ % https://x.ai/news/grok-3
        \cline{2-5}
                                         & \cellcolor{anthropicColor}{\textbf{Anthropic}} & Claude 3.7 Sonnet       & --                     & Yes \\ % https://www.anthropic.com/news/claude-3-7-sonnet
        \cline{2-5}
                                         & \cellcolor{alibabaColor}{\textbf{Alibaba}}     & QwQ-Max                 & MoE                   & Yes \\ % https://qwenlm.github.io/blog/qwq-max-preview/ (based on Qwen 2.5-Max, so MoE)
        \hline
        Mar 2025                         & \cellcolor{googleColor}{\textbf{Google}}       & Gemini 2.5              & --                     & Yes \\ % https://blog.google/technology/google-deepmind/gemini-model-thinking-updates-march-2025/
        \hline
        \multirow{2}{*}{Apr 2025}        & \cellcolor{metaColor}{\textbf{Meta}}           & LLaMA 4 109B/400B/2T    & MoE                   & -- \\ % https://ai.meta.com/blog/llama-4-multimodal-intelligence/
        \cline{2-5}
                                         & \cellcolor{openaiColor}{\textbf{OpenAI}}       & o3/o4-mini              & --                     & Yes \\ % https://cdn.openai.com/pdf/2221c875-02dc-4789-800b-e7758f3722c1/o3-and-o4-mini-system-card.pdf
        \hline
    \end{tabularx}
\end{table}

\pagebreak

\object{Major LLM Developers and Models}
The release of OpenAI's ChatGPT \cite{openai_introducing_2022} has innovated the research on \ac{llms}. Initially leveraging the GPT-3.5 architecture, OpenAI followed the trend of increasing secrecy with models like GPT-4 \cite{openai_gpt-4_2024} and its succesor GPT-4o \cite{openai_gpt-4o_2024}. While the \enquote{GPT} nomenclature suggests the continuation of the decoder-only architecture up to GPT-3 \cite{brown_language_2020}, many methodological details remain undisclosed. Recent models by OpenAI focused on \textit{reasoning} capabilities, as seen in the o1-mini \cite{openai_openai_2024}, o1 \cite{openai_openai_2024-1}, o3-mini \cite{openai_openai_2025}, and the subsequent o3 and o4-mini \cite{openai_openai_2025-1} releases.

Other major developers contributed to this competetive landscape. Meta started with a modified transformer, called LLaMA \cite{touvron_llama_2023}, and shifted to a decoder-only architecture in LLaMA 2 \cite{touvron_llama_2023-1} up to LLaMA 3.3\footnote{\url{https://github.com/meta-llama/llama-models/blob/main/models/llama3_3/MODEL_CARD.md}, accessed 05/09/2025. Note that auto-regressive means decoder-only}. With the release of the LLaMA 4 family \cite{meta_ai_llama_2025}, Meta switched to the \ac{moe} architecture, scaling its largest model to two trillion model parameters. Google followed a similar path, switching from its LaMDA-based Bard \cite{pichai_important_2023,thoppilan_lamda_2022} to the now well-known Gemini model series \cite{team_gemini_2024} by fine-tuning Bard. The family also started as decoder-only, while Gemini 1.5 marked a shift to an \ac{moe} architecture, though still described as decoder-based \cite{team_gemini_2024-1}. Beginning with Gemini 2.0 \cite{pichai_introducing_2024}, Google stopped publishing detailed architectural information. Gemini 2.5 \cite{kavukcuoglu_gemini_2025} further advanced their offerings by fully integrating reasoning capabilities, aligning with current state-of-the-art trends.

Besides the presented developers, Anthropic has remained private about architectural information of its Claude models, from Version 1 \cite{anthropic_introducing_2023}, Claude 2 \cite{anthropic_claude_2023}, through the Claude 3 family \cite{anthropic_introducing_2024-1}, Claude 3.5 Sonnet \cite{anthropic_introducing_2024} up to Claude 3.7 Sonnet \cite{anthropic_claude_2025}. The latter model also incorporated \textit{reasoning}, making it competitive with other currently leading models. Meanwhile, developers such as xAI, DeepSeek AI, and Alibaba have set up several \ac{moe} models. xAI's Grok-1 was openly an \ac{moe} model \cite{xai_open_2024}. However, for subsequent releases, Grok-2 \cite{xai_grok-2_2024} and Grok-3 \cite{xai_grok_2025}, architectural details were not provided, though Grok-3 is a competitive state-of-the-art model with \textit{reasoning}. DeepSeek AI has contributed powerful open-source \ac{moe} models like DeepSeek V3 \cite{deepseek-ai_deepseek-v3_2025} and its reasoning variant, R1 \cite{deepseek-ai_deepseek-r1_2025}. Similarly, the Alibaba Group has also contributed to the field with models like Qwen 2.5-Max \cite{team_qwen25-max_2025}, an \ac{moe} model. The more recent QwQ-Max \cite{team_thinkthink_2025} builds upon this architecture, incorporating reasoning to compete at the state-of-the-art.

Visible in Table \bfref{tab:llm_releases}, the current state-of-the-art for the given table began from o1 and o3-mini in December 2024 plus the release of DeepSeek V3. January 2025, DeepSeek R1 and Alibaba's Qwen 2.5-Max model without reasoning capability were released. February 2025 brought xAI's Grok-3, Anthropic's Claude 3.7 Sonnet and Alibaba's QwQ-Max. In March 2025, the state-of-the-art was reached by a major upgrade in Google Gemini, with the release of Gemini 2.5. Finally, April 2025 featured the LLaMA 4 family by Meta and OpenAI's o3 and o4-mini models. Based on this information pool, the models for the \ac{aqg} evaluation were selected, as to be seen in Section \bfref{sec:experiment-design}. 

\pagebreak

\object{Key \ac{llm} Development Trends} As illustrated in Table \bfref{tab:llm_releases}, a discernible trend is the decreasing transparency from developers regarding architectural designs of their \ac{llms}. While some models are known to be transformer-based (e.g.\ decoder-only) or employ an \ac{moe} architecture, and reasoning capabilities are often highlighted, many other details remain undisclosed due to the competitive nature of the field. It is also worth noting that there is a trend towards models with a higher number of parameters, a prevalence of decoder-only architectures from ChatGPT on, and an increasing adoption of \ac{moe} designs, a trend possibly popularized by GPT-4. Although the technical report for GPT-4 \cite{openai_gpt-4_2024} lacks architectural specifications, it is rumored to be an \ac{moe} model\footnote{\url{https://medium.com/@seanbetts/peering-inside-gpt-4-understanding-its-mixture-of-experts-moe-architecture-2a42eb8bdcb3}, accessed 05/09/2025}. Furthermore, a significant trend is the continuous integration of reasoning capabilities into \ac{llms}, with the most recent models being competitive in this area.

\subsection{Question Generation Research}
\label{sec:question-generation-research}

\ac{llms} enable the automation of the time-consuming and demanding procedure of question generation \cite{vu_chatgpt-based_2024}. As stated by \cite{naveed_comprehensive_2024}, the integration of \ac{llms} into educational settings offer an opportunity to improve the learning experience, support teachers and encourage the development of educational content. This results in freeing up time for educators to focus on teaching and interacting with students. Two certain directions of research are presented in the following:

\object{Prompt Engineering Approaches} Numerous studies have investigated the capability of \ac{llms} to generate educational questions, primarily aligned with cognitive levels, by solely refining and optimizing the prompts given to the models.

\cite{scaria_how_2024} addressed the problem of creating pedagogically effective questions, especially at higher skill levels, which is challenging, time-consuming, and requires expert-based evaluation. By this, researchers can understand the pedagogical aspects of \ac{aqg}, since automated metrics often analyze linguistic features. The study assessed the capability of \ac{llms}, including GPT-3.5 and GPT-4 and other models, to generate high-quality questions across different cognitive levels by using zero-shot prompting. The question evaluation was conducted by two annotators based on the categories adopted from \cite{horbach_linguistic_2020}, replacing the \textit{Rephrase} category with \textit{Bloom'sLevel}. This category assesses the adherence to the targeted cognitive level. The Inter-annotator agreement was measured using Cohen's Kappa \cite{cohen_coefficient_1960}, since there are two annotators to be validated. Their findings revealed that over $91\%$ of the generated questions were relevant and high-quality, where models such as GPT-4 and GPT-3.5 were notably effective, especially in adhering to the desired Bloom level. However, limitations were identified by the authors, e.g.\ the crucial human evaluation can be subjective. Also, most models were unable to generate high-quality questions at the \textit{Apply} and \textit{Create} levels, and the study lacked prompt diversity, as the same zero-shot prompt was used for all models.

\pagebreak

\cite{maity_can_2025} have described the gap that frameworks like Bloom's Taxonomy are often not fully integrated into the systematic assessment of generated questions, or they have not employed comprehensive methods to determine the completeness and relevance of the questions. Their research explored the potential of \ac{llms} -- such as several GPT and LLaMA models plus Gemini Pro -- for \ac{aqg} by comparing zero-shot prompting against eight-shot prompting. By this, they generated a complete set of questions aligned with Bloom's Taxonomy using school textbooks. Using the same evaluation metrics as \cite{elkins_how_2023}, the teachers-based evaluation indicated that eight-shot prompting improved the quality and Bloom adherence of educational questions. GPT-4 was identified as the best performing model under both prompting conditions. As future work, the authors suggested exploring 16- and 32-shot prompting, which could further enhance the quality of generated questions. However, they have not addressed the subjective nature of human evaluation, because 28 teachers were involved in the evaluation process, which could lead to a more diverse assessment of the generated questions.

The authors of \cite{al_faraby_analysis_2024} also stated that the quality and pedagogical effectiveness of \ac{llm}-generated questions have not been thoroughly evaluated. Furthermore, \ac{llm}-based question generation, which aims to meet educational objectives, is unexplored. The assessment of these questions plus the comparison with expert-generated ones is the focus of their study, primarily using ChatGPT. While disregarding frameworks such as Bloom's Taxonomy, the study focused on textbook-based questions and different prompting strategies: zero-shot, enhanced zero-shot, few-shot, and few-shot with \ac{cot}, where each strategy was tested with 6 variations. Human evaluators -- experts and crowdsourced workers -- were employed to assess the generated questions. The results revealed that experts preferred human-generated questions from textbooks, while the crowdsourced evaluators have shown comparable preferences for both human- and \ac{llm}-generated questions. However, the authors also highlighted the issue of hallucination, which remains a challenge in \ac{aqg}. They suggested that future work should focus on integrating systems to verify questions against source texts. Moreover, the study had a small -- but unknown -- number of annotators. They suggested a broader range of \ac{llms} and a higher variety of questions to be included in future work.

Traditional \ac{aqg} systems struggle to generate questions at higher cognitive levels, and they often lack control over question complexity and difficulty \cite{blobstein_angel_2023}. Prior work has focused on incorporating pre-formulated answers into the question generation process. Addressing this, the authors developed Angel, a tool that utilizes GPT-3.5 for generating questions and answers based on learning material. The system employs instructional prompts for \ac{aqg} and evaluation -- including few-shot prompting, descriptions for the Bloom levels, and textbook content to produce questions categorized by difficulty (easy, medium, hard). On the one hand, the evaluation of the generated questions was conducted using GPT-3.5 again as an evaluator with a \ac{cot} strategy, focusing on metrics such as \textit{Clarity, Relatedness, Importance,} and \textit{Answerability}. On the other hand, human evaluation was also performed, coupled with GPT-3.5 for assessing the resulted Bloom levels of the questions. While the LLM-based evaluation indicated that Angel generated more questions at higher cognitive levels -- in contrast to the contestant systems -- the human evaluation underscored that the majority of the questions were still classified as \textit{Remembering}. The authors concluded that while Angel demonstrated potential for \ac{aqg}, further work and human evaluation are needed to analyze the results more accurately.

\object{Hybrid Architectures} Researchers have also developed more complex systems that integrate \ac{llms} with other components in a coordinated fashion. These hybrid architectures often aim to leverage the strengths of different techniques to achieve more robust and diverse question generation.

In contrast to the previous studies, \cite{doughty_comparative_2024} addressed the gap of generating \ac{mcqs} in programming education, which is linked with much effort. \ac{llms} were not thoroughly explored for this purpose, since evaluating systems which generate \ac{mcqs} is limited. The authors conducted a comparative study of \ac{mcqs} generated by GPT-4 versus those created by humans. In addition to the \ac{llm}, few-shot prompting, course context and learning objectives aligned with Bloom's Taxonomy levels -- realized through fine-tuned BERT classifier -- were used for \ac{aqg}. The resulted questions were comparable to human-generated ones in terms of clarity and learning objective coverage. Nevertheless, limitations included potential bias from human evaluators, the focus on programming questions, and the need to consider the alignment between the question and the learning content, as depicted similarly in \cite{al_faraby_analysis_2024} above.

Focusing on personalized learning, \cite{hang_mcqgen_2024} described that both efficient and effective generation of \ac{mcqs} is essential in modern teaching, since the manual generation is depicted as resource-intensive and time-consuming. Their system, MCQGen, utilizes GPT-4, \ac{rag}, and advanced prompting techniques to automate the creation of personalized \ac{mcqs} while accessing an external knowledge base. The self-refinement and \ac{cot} prompting were employed to enhance the quality of the generated questions. Assessing the questions based on \textit{Grammar, Answerability, Diversity, Complexity,} and \textit{Relevance}, the outcomes between the three \ac{llm}-based evaluators and three lecturers were mixed. In this work, especially the \textit{Complexity} and \textit{Diversity} metrics were lacking in both evaluation methods, suggesting that prompt optimization and expanding the knowledge base could address these limitations. In addition to that, the usage of only one \ac{llm} coupled with the focus on computer science questions in English language were also limitations of this work.

\subsection{Research Gaps Addressed}

Based on the conducted literature review, presented in the previous sections of the thesis, there are several research gaps this work aims to address: 

\object{Outdated Models} Current studies rely on \ac{llms} from 2023 or earlier, such as GPT-4, GPT-3.5, and LLaMA2 \cite{scaria_how_2024, maity_can_2025, al_faraby_analysis_2024}. These evaluations lack in capturing the capabilities of recent models that incorporate advanced \textit{reasoning} mechanisms to think before answering. To address this, the thesis investigates four current state-of-the-art models, as detailed in Section \bfref{sec:experiment-design}. This ensures that the evaluation of \ac{aqg} capabilities is conducted using the most recent technologies, resulting in a more up-to-date understanding of \ac{llms}' potential in educational contexts.

\pagebreak

\object{Limited Content Adherence Evaluation} A persistent challenge in \ac{aqg} is the problem of hallucination, which is unacceptable in education \cite{yang_heuristic_2024}. \ac{llms} can generate misleading or irrelevant information which are not directly supported by e.g.\ the source text \cite{al_faraby_analysis_2024,duong-trung_bloomllm_2024}. While some studies mention this issue, the content adherence between a provided context and the generated questions still remains underexplored, as stated in the limitations of e.g.\ \cite{doughty_comparative_2024,al_faraby_analysis_2024}. This thesis investigates this problem in RQ1 -- conducting both quantitative and qualitative analysis. This kind of evaluation is crucial for generating reliable educational questions that can be answered accurately using the source material, as diverse materials were used in e.g.\ \cite{scaria_how_2024,maity_can_2025,blobstein_angel_2023}. Highlighted in related research \cite{biancini_multiple-choice_2024}, it is important to inject knowledge directly into prompts to contrast hallucinations and give educators control over the source text.

\object{Bloom's Taxonomy Integration and qualitative Evaluation} While several studies have incorporated Bloom's Taxonomy into \ac{aqg} evaluation \cite{scaria_how_2024, maity_can_2025, blobstein_angel_2023}, most research lacks comprehensive analysis of how different question formats perform across various cognitive levels. Existing work primarily focuses on adherence to specified Bloom levels without examining the relationship between question types (e.g.\ Multiple-Choice vs. Open-Ended) and their pedagogical effectiveness at different cognitive complexities. In addition to that, the \ac{llms} used struggle with higher cognitive levels \cite{blobstein_angel_2023}, such as \textit{Apply} and \textit{Create} \cite{scaria_how_2024}. Since these models are far from the state-of-the-art, there is an opportunity to explore how the latest \ac{llms} perform in generating questions at all six Bloom levels, while several publications did not include Bloom's Taxonomy into their work \cite{al_faraby_analysis_2024,vu_chatgpt-based_2024,hang_mcqgen_2024}. 

Furthermore, many studies rely heavily on automated metrics or limited human evaluation \cite{zhuge_twinstar_2025,blobstein_angel_2023,li_planning_2024} -- often with few annotators --, which might not capture the full pedagogical value of generated questions. This thesis addresses this mixed gap primarily through RQ2 by conducting a comprehensive analysis of question format performance across all six Bloom levels, employing expert-based evaluation with established inter-rater agreement measures. The evaluation framework incorporates rubrics that assess not only Bloom level adherence but also qualitative criteria, providing a proper understanding of \ac{llm} capabilities in \ac{aqg}. The rubrics for both experiments are proposed in Section \bfref{sec:experiment-design} and are based on the work \cite{mi_comparative_2024}.

This set of research gaps highlights the need for a comprehensive evaluation that can evaluate modern \ac{llms} in \ac{aqg} using current methods. The proposed framework in the section below aims to bridge the research gaps and provide a foundation for future research in \ac{aqg} and \ac{llms}.